\section{Log Management System}

\subsection{Definisi}
Pemantauan server melalui aktivitas \textit{logging} merupakan elemen krusial bagi administrator sistem untuk melacak aktivitas pengguna, menyeimbangkan beban (\textit{load balancing}), hingga mendeteksi serangan seperti DDoS. Namun, sistem \textit{logging} konvensional seringkali membebani sumber daya karena biaya yang tinggi dalam hal penyimpanan, pengumpulan data, serta proses pencarian dan analisis.

Implementasi manajemen log yang efektif mengikuti siklus hidup yang terstruktur untuk memastikan kegunaan data (Chuvakin et al., 2012):
\begin{enumerate}
    \item \textbf{Log Generation}: Proses pembuatan entitas data oleh aplikasi, server, atau perangkat jaringan.
    \item \textbf{Log Analysis}: Penggunaan algoritma untuk mengidentifikasi pola anomali atau ancaman keamanan.
    \item \textbf{Log Storage \& Retention}: Penyimpanan data dalam jangka waktu tertentu sesuai kebijakan retensi (misal: log akses rekam medis wajib disimpan minimal 5 tahun).
    \item \textbf{Log Disposal}: Penghapusan data log secara aman setelah masa retensi berakhir untuk menjaga efisiensi penyimpanan.
\end{enumerate}

\subsection{Contoh Implementasi}
Nguyen et al. \cite{18}, mengusulkan sebuah sistem bernama eLMS (\textit{efficient Log Management System}). Sistem ini dirancang untuk dapat mengumpulkan log dari berbagai server secara terukur (\textit{scalable}), terindeks dengan baik, dan dapat dianalisis dengan cepat. Secara teknis, eLMS dikembangkan berbasis core \textit{open-source} ELK (ElasticSearch, LogStash, Kibana), namun dengan modifikasi signifikan pada metode pengiriman datanya. Alih-alih menggunakan teknik \textit{upload} standar seperti pada LogStash, sistem ini menerapkan teknik antrean (\textit{queue}) dan \textit{streaming}. Hasil pengujian menunjukkan bahwa pendekatan ini mampu meningkatkan kecepatan pemrosesan hingga sepuluh kali lipat dibandingkan sistem ELK standar, menjadikannya solusi yang lebih efisien untuk manajemen log skala besar.